\section{Verification plan}
\label{sec:verify}

The architecture of \S\ref{sec:arch} is verified against the golden C model,
and the golden C model is verified against the measurements. This section
states both halves and what each one can and cannot establish.

\subsection{The golden model as an oracle}

The golden model is a deterministic, transaction-level C++17 model of one
endpoint, with the caller owning the clock. It exposes a C-linkage facade over
plain fixed-width structures with picosecond timestamps and no exceptions
crossing the boundary: construct an endpoint from a profile, post work
requests, ring a doorbell, deliver a received packet, deliver a fabric event,
advance to a timestamp, ask for the next internally scheduled time, poll
completions, drain emitted packets, read the counter facade, write a
transaction trace, destroy.

Two properties make it usable as an oracle rather than as a second opinion.
First, determinism is a contract: the same stimulus sequence and the same
profile produce byte-identical traces, so a trace recorded from the simulator
is an expected-result file, and a divergence localizes to the first differing
timestamp or counter rather than to a curve that looks different. Second, the
next-event query lets a testbench step the model exactly, so the comparison is
not confounded by the testbench's own time quantisation.

\subsection{DPI-C in a UVM environment}

The design under test and the golden model run side by side under one
stimulus. The C facade is imported through the direct programming interface,
and the environment is arranged as follows.

\textbf{Stimulus} has two sources, both randomized under constraints derived
from the campaign's own sweep dimensions: work requests (opcode, message size,
gather-entry count, signalling, destination, posting batch, queue depth,
inter-burst gap) driven into the host interface as doorbells and host-memory
contents, and packets (data, acknowledgement, negative acknowledgement,
notification, pause) driven into the media-access interface. The gap parameter
matters more than it looks: it is the axis along which \anom{04} lives, and a
constrained-random generator that only ever saturates will never enter the
lossless regime that serving traffic actually occupies.

\textbf{The scoreboard} compares three things: the sequence of emitted packets,
field by field with the mutable fields masked exactly as the invariant check
masks them; the sequence of completion-queue entries with their work-request
identity, status and byte count; and, at every snapshot point, the full counter
set. Timestamps are compared within a per-stage band rather than exactly,
because the model's stage services are fitted rather than cycle accurate, and
the band is part of each block's specification. Counters are compared exactly.

\textbf{Coverage and assertions come from the anomaly table.} Every row of
Table~\ref{tab:anomaly} becomes both a coverage bin and a check.
A row whose kind is \textit{emergent} is a check on the mechanism: the
condition in its trigger column must be reached (coverage) and the effect must
land inside its registered band (assertion), with the band taken from the
model's own validation contract, for example a 15\,\% band on the depth-1
goodput curve, a 20\,\% band on the depth ratio, a 78 to 92\Gbps{} window on
the saturated single-queue-pair equilibrium, three percentage points on the
multi-queue-pair ceiling, two percentage points around the 5.6\,\% maximum
transmission unit tax, 25\,\% on the incast tax identity with two points on the
fair-share split, 30\,\% around 283 notifications per second per congested
queue pair, and exactly zero notifications for a lone paced flow below
saturation. A row whose kind is \textit{injected} is a check that the rule
fired, not that a mechanism produced it, and the report says so. A row whose
kind is \textit{fabric} is not a check on the endpoint at all and is
implemented in the fabric model. A row whose kind is \textit{counter} is a
check on the facade with an explicit assertion that the datapath did
\emph{not} change. A row whose kind is \textit{tool} is reproduced by nothing
and exists so that nobody re-derives a number from an instrument artifact.

\textbf{Trace replay} closes the loop with the real device. The campaign's
curated result rows are golden vectors: the depth-1 and depth-1024 message-size
curves, the gap sweep at 8 and 64\,KiB, the multi-queue-pair ceiling, the
maximum-transmission-unit pair, the bidirectional pair, the incast rows and the
fan-out control. Each is registered in an expectations file before the run and
reports as one row with a band. Three fatal guards void a run rather than
costing it a point: deterministic replay must be identical across two runs,
the one-authority conservation check must hold, and counters must be monotone.

\subsection{The full chain}

Above the endpoint sits a packet-level fabric model, which owns links, switch
queues, marking, drops and pause, while the endpoint owns work-queue elements,
contexts, transport state and rate. The seam between them carries packet
attempts and fabric events and nothing else: no queue, context or completion
object crosses it. For this deployment the fabric is configured as the
measurements found it, that is, one non-blocking switch, four 100\Gbps{} ports,
515\nS{} pipes, a 5.2\,MB tail-drop egress buffer per port, no marking, no
priority flow control and no propagated pause. Configuring the switch to be
silent is what forces the endpoint to notice congestion by itself, which is the
whole point of \anom{16}, and it is why the fabric model's random-early-detection
policy is disabled for this cluster rather than tuned.

The chain matters because two of the open clauses are only visible in it. The
congestion-control loop settles wherever the ingress meter's drain rate is; at
the drain rate fitted from the lone-flow anchors the loop settles below the
switch's egress rate, the switch buffer never fills, there is no loss to
amplify, and the incast tax comes out near zero instead of at the measured
26.9\,\%. The incast measurement is a second anchor for the same latent and it
wants a higher value. One drain rate cannot satisfy both, which is the same
shape as the duplex-versus-simplex incompatibility already registered against
that block, and it points at the same missing stall mechanism.

\subsection{What the FPGA can show that the silicon cannot}
\label{sec:verify:fpga}

The measured device publishes about a hundred counters, non-atomically, cached
for 10\mS{} by the driver, several of them meaning something other than their
names. The implementation publishes, through the monitoring ring and at one
coherent snapshot instant, the occupancy and high-water mark of every internal
queue, the credit state of every arbiter, the per-stage service counts that the
fixed-offset cost $T_{\text{eff}}$ currently lumps into one number, and the
reaction point's rate and $\alpha$ per queue pair.

Three campaign questions become answerable there and only there. The split of
$T_{\text{eff}}$ across publication, observation, fetch, context readiness and
admission, which is \evcal{} today, becomes five measured numbers. The
\anom{03} discard ledger, where the counted discards are far too few to explain
the goodput deficit, becomes an occupancy trace of the ingress meter at the
operating point, which either shows the buffer filling or does not. And
\anom{17}, the stall burst whose mechanism is unknown, becomes a time series of
ingress occupancy against arrival rate during a burst, with the reverse-traffic
dependence as a controlled input. None of these requires the vendor's device to
cooperate; they require an implementation whose observation path was designed
in from the start, which is why the monitoring network on chip is in the
architecture rather than in a debug appendix.
